\documentclass[10pt,a4paper]{amsart}
\usepackage{amsmath,amssymb,amsthm,graphicx}
% Retained preamble for the Stacks Project chapter snapshot in this directory.
% The chapter's own preamble.tex is a copy of the shared Stacks Project preamble; the block
% below is sliced MECHANICALLY, line for line, from it. Nothing was added, removed or
% redefined; the wrapper class, the xr-hyper/hyperref wiring and the externaldocument list
% are omitted because the assembler supplies the document class.
%
% lines taken from preamble.tex: 3, 5, 9, 11, 12, 156, 157, 158, ... (51 lines total)
%
% package coverage restored from the chapter's own preamble (2026-10-13): xy, xy-2cell

\usepackage[all]{xy}
\xyoption{2cell}
\UseAllTwocells

\newenvironment{reference}{\comment}{\endcomment}
\newenvironment{slogan}{\comment}{\endcomment}
\newenvironment{history}{\comment}{\endcomment}
\theoremstyle{plain}
\newtheorem{theorem}[subsection]{Theorem}
\newtheorem{proposition}[subsection]{Proposition}
\newtheorem{lemma}[subsection]{Lemma}
\theoremstyle{definition}
\newtheorem{definition}[subsection]{Definition}
\newtheorem{example}[subsection]{Example}
\newtheorem{exercise}[subsection]{Exercise}
\newtheorem{situation}[subsection]{Situation}
\theoremstyle{remark}
\newtheorem{remark}[subsection]{Remark}
\newtheorem{remarks}[subsection]{Remarks}
\numberwithin{equation}{subsection}
\def\lim{\mathop{\mathrm{lim}}\nolimits}
\def\colim{\mathop{\mathrm{colim}}\nolimits}
\def\Spec{\mathop{\mathrm{Spec}}}
\def\Hom{\mathop{\mathrm{Hom}}\nolimits}
\def\Ext{\mathop{\mathrm{Ext}}\nolimits}
\def\SheafHom{\mathop{\mathcal{H}\!\mathit{om}}\nolimits}
\def\SheafExt{\mathop{\mathcal{E}\!\mathit{xt}}\nolimits}
\def\Sch{\mathit{Sch}}
\def\Mor{\mathop{\mathrm{Mor}}\nolimits}
\def\Ob{\mathop{\mathrm{Ob}}\nolimits}
\def\Sh{\mathop{\mathit{Sh}}\nolimits}
\def\NL{\mathop{N\!L}\nolimits}
\def\CH{\mathop{\mathrm{CH}}\nolimits}
\def\proetale{{pro\text{-}\acute{e}tale}}
\def\etale{{\acute{e}tale}}
\def\QCoh{\mathit{QCoh}}
\def\Ker{\mathop{\mathrm{Ker}}}
\def\Im{\mathop{\mathrm{Im}}}
\def\Coker{\mathop{\mathrm{Coker}}}
\def\Coim{\mathop{\mathrm{Coim}}}
\DeclareMathSymbol{\boxtimes}{\mathbin}{AMSa}{"02}
\def\QCohstack{\mathcal{QC}\!\mathit{oh}}
\def\Cohstack{\mathcal{C}\!\mathit{oh}}
\def\Spacesstack{\mathcal{S}\!\mathit{paces}}
\def\Quotfunctor{\mathrm{Quot}}
\def\Hilbfunctor{\mathrm{Hilb}}
\def\Curvesstack{\mathcal{C}\!\mathit{urves}}
\def\Polarizedstack{\mathcal{P}\!\mathit{olarized}}
\def\Complexesstack{\mathcal{C}\!\mathit{omplexes}}
\def\Pic{\mathop{\mathrm{Pic}}\nolimits}
\def\Picardstack{\mathcal{P}\!\mathit{ic}}
\def\Picardfunctor{\mathrm{Pic}}
\def\Deformationcategory{\mathcal{D}\!\mathit{ef}}
\title{The Chinese remainder theorem for pairwise comaximal ideals}
\author{The Stacks Project Authors}
\date{stacks-project:algebra.tex:local-snapshot-20261012; GNU FDL 1.2 or later}
\begin{document}
\maketitle
\noindent\textit{Editorial note.} All delivered mathematical text is the author's own, copied from the pinned Stacks Project chapter file algebra.tex; no statement or proof step was rewritten, reordered or removed.\par\medskip
\noindent\textit{Editorial note.} Preamble provenance, as required for this source: the frozen snapshot stages the chapter's own preamble file beside it (retained-preamble.tex, 63 lines, 8116 bytes, the same shared Stacks-wide preamble the published chow and etale-cohomology cards use); the first line of the pinned chapter file is \textbackslash{}input\{preamble\}, so no preamble text lives inside the chapter itself. The card takes that file's lines 8-18 (the verbatim package and the three \textbackslash{}comment-based environments reference/slogan/history at lines 9-12, section 13.2, together with the Xy-pic setup) and its lines 154 to the runtime end of the file (the native theorem declarations and mathematical macros). Not one character was changed, no package was added, no macro was defined, and the excerpt serves only as preamble, never as a span coordinate.\par\medskip
\noindent\textit{Editorial note.} The same-chapter pointers are located at their \textbackslash{}label lines inside the pinned chapter file; the cross-chapter pointers are located in the project's own tag list staged in the snapshot (sources/stacks-chow/tags.txt), at the verbatim lines carrying their canonical tags.\par\medskip
\noindent\textit{Editorial note.} The counted spans contain no \textbackslash{}ref and no \textbackslash{}cite, so the card carries no navigation stub and no bibliography.\par\medskip
\noindent\textit{Allowed standard prerequisites (editorial).}\begin{itemize}
\item The product of ideals is contained in their intersection. Used at the start of the counted proof (line 2850). Gives one of the two inclusions for the intersection identity.
\item Comaximality to produce a partition of unity. Used in the second half of the counted proof. Produces elements whose sums give the identity modulo each ideal, which yields surjectivity.
\end{itemize}\medskip
\begin{lemma}[Chinese remainder]
\label{lemma-chinese-remainder}
Let $R$ be a ring.
\begin{enumerate}
\item If $I_1, \ldots, I_r$ are ideals such that $I_a + I_b = R$
when $a \not = b$, then $I_1 \cap \ldots \cap I_r =
I_1I_2\ldots I_r$ and $R/(I_1I_2\ldots I_r)
\cong R/I_1 \times \ldots \times R/I_r$.
\item If $\mathfrak m_1, \ldots, \mathfrak m_r$ are pairwise distinct maximal
ideals then $\mathfrak m_a + \mathfrak m_b = R$ for $a \not = b$ and the
above applies.
\end{enumerate}
\end{lemma}

\begin{proof}
Let us first prove $I_1 \cap \ldots \cap I_r = I_1 \ldots I_r$
as this will also imply the injectivity of the induced ring
homomorphism $R/(I_1 \ldots I_r) \rightarrow R/I_1 \times \ldots \times R/I_r$.
The inclusion $I_1 \cap \ldots \cap I_r \supset I_1 \ldots I_r$ is always
fulfilled since ideals are closed under multiplication with arbitrary ring
elements. To prove the other inclusion, we claim that the ideals
$$
I_1 \ldots \hat I_i \ldots I_r,\quad i = 1, \ldots, r
$$
generate the ring $R$. We prove this by induction on $r$. It holds when
$r = 2$. If $r > 2$, then we see that $R$ is the sum of the ideals
$I_1 \ldots \hat I_i \ldots I_{r - 1}$, $i = 1, \ldots, r - 1$.
Hence $I_r$ is the sum of the ideals
$I_1 \ldots \hat I_i \ldots I_r$, $i = 1, \ldots, r - 1$.
Applying the same argument with the reverse ordering on the ideals
we see that $I_1$ is the sum of the ideals
$I_1 \ldots \hat I_i \ldots I_r$, $i = 2, \ldots, r$.
Since $R = I_1 + I_r$ by assumption we see that $R$ is the sum of the
ideals displayed above. Therefore we can find elements
$a_i \in I_1 \ldots \hat I_i \ldots I_r$
such that their sum is one. Multiplying this equation by an element
of $I_1 \cap \ldots \cap I_r$ gives the other inclusion.
It remains to show that the canonical map
$R/(I_1 \ldots I_r) \rightarrow R/I_1 \times \ldots \times R/I_r$
is surjective. For this, consider its action on the equation
$1 = \sum_{i=1}^r a_i$ we derived above. On the one hand, a
ring morphism sends 1 to 1 and on the other hand, the image of any
$a_i$ is zero in $R/I_j$ for $j \neq i$. Therefore, the image of $a_i$
in $R/I_i$ is the identity. So given any element
$(\bar{b_1}, \ldots, \bar{b_r}) \in R/I_1 \times \ldots \times R/I_r$,
the element $\sum_{i=1}^r a_i \cdot b_i$ is an inverse image in $R$.

\medskip\noindent
To see (2), by the very definition of being distinct maximal ideals, we have
$\mathfrak{m}_a + \mathfrak{m}_b = R$ for $a \neq b$ and so the above applies.
\end{proof}

\end{document}
